\documentclass[conference]{IEEEtran}
\usepackage[utf8]{inputenc}
\usepackage[T1]{fontenc}
\usepackage[spanish,es-tabla]{babel}
\usepackage{graphicx}
\usepackage{xcolor}
\usepackage{array}
\usepackage{booktabs}
\usepackage{multirow}
\usepackage{enumitem}
\usepackage{tikz}
\usetikzlibrary{positioning,arrows.meta,shapes.geometric,fit,backgrounds}
\usepackage[most]{tcolorbox}
\usepackage[hidelinks]{hyperref}
\usepackage{url}

\definecolor{ink}{HTML}{1E293B}
\definecolor{soft}{HTML}{475569}
\definecolor{accent}{HTML}{0891B2}
\definecolor{accent2}{HTML}{7C3AED}
\definecolor{ok}{HTML}{16A34A}
\definecolor{warn}{HTML}{EA580C}
\definecolor{bad}{HTML}{DC2626}
\definecolor{paper}{HTML}{F8FAFC}

\newcolumntype{L}[1]{>{\raggedright\arraybackslash}p{#1}}
\newcolumntype{C}[1]{>{\centering\arraybackslash}p{#1}}
\newcommand{\pkg}[1]{\texttt{#1}}
\newcommand{\attr}[1]{{\scriptsize\color{soft}#1}}

\newtcolorbox{nota}[1][]{
  colback=paper, colframe=accent, boxrule=0.6pt, arc=2pt,
  left=6pt, right=6pt, top=4pt, bottom=4pt,
  fonttitle=\bfseries\footnotesize, title={#1}
}
\newtcolorbox{alerta}[1][]{
  colback=red!4, colframe=bad, boxrule=0.6pt, arc=2pt,
  left=6pt, right=6pt, top=4pt, bottom=4pt,
  fonttitle=\bfseries\footnotesize, title={#1}
}

\begin{document}

\title{EnRuta: Plataforma de Log\'istica de \'Ultima Milla Basada en Microservicios Dirigidos por Eventos}
\author{\IEEEauthorblockN{Kevin Rachez, Ronald Vargas, Santiago Delgadillo}
\IEEEauthorblockA{Programa de Ingenier\'ia de Sistemas, Universidad de Cundinamarca\\
Fusagasug\'a, Cundinamarca, Colombia}
\and
\IEEEauthorblockN{Prof. Frank Ortengon}
\IEEEauthorblockA{Docente asesor, Universidad de Cundinamarca\\
Fusagasug\'a, Colombia}}
\maketitle

\markboth{EnRuta: Log\'istica de \'ultima milla con microservicios --- Rachez, Vargas, Delgadillo, Ortengon}{}

\begin{abstract}
EnRuta es una plataforma de log\'istica de \'ultima milla construida sobre siete microservicios
desplegados con Docker y verificados de extremo a extremo sobre el casco urbano de Fusagasug\'a
(Colombia). Resuelve tres problemas que un CRUD no resuelve: asignar un repartidor sin duplicados
bajo contenci\'on concurrente, propagar los cambios de estado sin perder eventos y rastrear el
env\'io en tiempo real con GPS. La soluci\'on combina asignaci\'on con triple barrera anti-doble,
patr\'on \emph{outbox} con relay transaccional, saga coreografiada, geocercas con hist\'eresis e
idempotencia extremo a extremo. El sistema se valid\'o con 367 pruebas automatizadas (97,7\,\% de
cobertura de l\'ineas), prueba de concurrencia de 50 hilos con un \'unico ganador, ciclo completo
\texttt{CREATED} $\rightarrow$ \texttt{DELIVERED} en $\sim$2 min con correo real entregado en un
intento, ETA real OSRM de 2\,529 m / 273 s con cach\'e en la segunda llamada y plan B de 43 ms, y
observabilidad con 8/8 objetivos Prometheus en estado UP.
\end{abstract}

\begin{IEEEkeywords}
microservicios, log\'istica de \'ultima milla, Apache Kafka, patr\'on outbox, saga, idempotencia,
geocercas, Spring Boot, Nominatim, OSRM, Fusagasug\'a
\end{IEEEkeywords}

\section{Introducción}
% ============================================================================

\subsection{Contexto}

La \textbf{última milla} es el tramo final de una entrega: desde que el paquete sale de un punto
de recogida hasta que llega a manos del cliente. Aunque parece un problema simple, concentra las
mayores dificultades de una operación logística: asignar el repartidor adecuado, seguirlo en
tiempo real, avisar al cliente y no perder ningún evento cuando algo falla.

\textbf{EnRuta} es una plataforma que resuelve ese problema para una ciudad intermedia
(Fusagasugá, Cundinamarca) y que se ejecuta íntegramente en un computador portátil. Un cliente
pide un envío desde el navegador; un repartidor lo acepta desde su teléfono; el sistema lo asigna,
lo rastrea con GPS real y envía un correo al entregarlo.

\subsection{Problema}

Construir el flujo completo exige resolver cuatro problemas que un CRUD no resuelve:

\begin{enumerate}[leftmargin=1.2cm]
  \item \textbf{Asignar sin duplicar.} Dos repartidores pueden aceptar la misma oferta en el
        mismo milisegundo; el pedido no puede quedar asignado a dos personas.
  \item \textbf{No perder eventos.} El cambio de estado y el aviso a los demás servicios deben
        ocurrir juntos: no puede existir un envío sin evento ni un evento sin envío.
  \item \textbf{Seguir en vivo.} Las posiciones GPS llegan cientos de veces por minuto; el mapa
        debe actualizarse sin recargar y sin bloquear al cliente.
  \item \textbf{Sobrevivir a los fallos.} El correo puede caer, el motor de rutas puede no
        responder y un repartidor puede abandonar el pedido a mitad de camino.
\end{enumerate}

\subsection{Solución propuesta}

Una arquitectura de \textbf{siete microservicios} donde cada uno posee su propio dato y se
comunica con los demás mediante eventos (Kafka) o llamadas HTTP puntuales. La tabla
\ref{tab:servicios} resume los servicios y la figura \ref{fig:arquitectura} su disposición.

\begin{table*}[t]
\centering
\caption{Servicios del sistema}
\label{tab:servicios}
\small
\begin{tabular}{@{}L{2.6cm}L{1.2cm}L{5.6cm}L{4.2cm}@{}}
\toprule
\textbf{Servicio} & \textbf{Puerto} & \textbf{Responsabilidad} & \textbf{Dato que posee}\\
\midrule
\pkg{api-gateway} & 18080 & valida JWT, enruta y limita peticiones & —\\
\pkg{auth-service} & 8081 & usuarios, login, refresh y JWKS & \pkg{authdb}\\
\pkg{shipment-service} & 8082 & ciclo de vida del envío e historial & \pkg{shipmentdb}\\
\pkg{dispatch-service} & 8083 & ofertas y asignación de repartidores & \pkg{dispatchdb} + Redis\\
\pkg{tracking-service} & 8084 & GPS, geocercas y WebSocket & Redis (sin BD)\\
\pkg{routing-eta-service} & 8085 & rutas y ETA con OSRM & caché en Redis\\
\pkg{notification-service} & 8086 & correos con reintentos y DLQ & \pkg{notifdb}\\
\bottomrule
\end{tabular}
\end{table*}

\begin{figure*}[t]
\centering
\resizebox{\textwidth}{!}{%
\begin{tikzpicture}[
  node distance=0.9cm and 1.1cm,
  box/.style={draw, rounded corners, minimum width=2.6cm, minimum height=0.9cm,
              align=center, font=\small, fill=white},
  infra/.style={draw, rounded corners, minimum width=2.4cm, minimum height=0.7cm,
                align=center, font=\scriptsize, fill=paper},
  arrow/.style={-{Latex}, thick, color=soft}
]
  \node[box, fill=accent!8, draw=accent] (web) {Cliente\\ \scriptsize web / PWA};
  \node[box, fill=accent2!8, draw=accent2, right=of web] (gw) {api-gateway};
  \node[box, fill=ok!8, draw=ok, right=of gw] (ship) {shipment};
  \node[box, fill=warn!8, draw=warn, below=of ship] (disp) {dispatch};
  \node[box, fill=accent!8, draw=accent, below=of disp] (track) {tracking};
  \node[box, fill=bad!8, draw=bad, below=of track] (notif) {notification};
  \node[box, fill=ok!8, draw=ok, above=of ship] (auth) {auth};
  \node[box, fill=accent2!8, draw=accent2, above=of auth] (routing) {routing-eta};

  \draw[arrow] (web) -- (gw);
  \draw[arrow] (gw.east) -- ++(0.4,0) |- (ship.west);
  \draw[arrow] (gw.east) -- ++(0.4,0) |- (disp.west);
  \draw[arrow] (gw.east) -- ++(0.4,0) |- (track.west);
  \draw[arrow] (gw.east) -- ++(0.4,0) |- (notif.west);
  \draw[arrow] (gw.east) -- ++(0.4,0) |- (auth.west);
  \draw[arrow] (gw.east) -- ++(0.4,0) |- (routing.west);

  \node[infra, below=1.4cm of notif, xshift=-3cm] (kafka) {Kafka};
  \node[infra, right=0.4cm of kafka] (pg) {PostgreSQL};
  \node[infra, right=0.4cm of pg] (redis) {Redis};
  \node[infra, right=0.4cm of redis] (osrm) {OSRM};
  \node[infra, right=0.4cm of osrm] (mail) {MailHog};
  \node[draw=soft, dashed, rounded corners, fit=(kafka)(pg)(redis)(osrm)(mail), inner sep=6pt] {};
\end{tikzpicture}
}
\caption{Arquitectura de contenedores. Todos los servicios comparten infraestructura, pero cada
uno tiene su propio esquema de datos.}
\label{fig:arquitectura}
\end{figure*}

\subsection{Alcance}

El sistema cubre el ciclo completo: creación del envío, asignación, recogida, tránsito, entrega y
notificación. Quedan fuera del alcance: pagos reales, multi-empresa, aplicación móvil nativa y
facturación fiscal. El capítulo \ref{ch:futuro} detalla las líneas de trabajo futuras.

% ============================================================================
\clearpage
\section{Objetivos}
% ============================================================================

\subsection{Objetivo general}

Diseñar e implementar una plataforma de logística de última milla basada en microservicios, con
asignación de repartidores consistente, rastreo GPS en tiempo real y notificaciones confiables.

\subsection{Objetivos específicos}

\begin{enumerate}[leftmargin=1.2cm]
  \item Modelar el dominio y las bases de datos con independencia por servicio.
  \item Garantizar la consistencia del estado mediante el patrón \emph{outbox} y una saga
        coreografiada.
  \item Evitar la doble asignación con tres mecanismos independientes.
  \item Implementar un pipeline de ingesta GPS con geocercas e histéresis y un canal WebSocket.
  \item Proveer idempotencia en las escrituras y deduplicación en los consumidores.
  \item Alcanzar una cobertura de pruebas de al menos el 70\,\% de las líneas.
  \item Documentar las decisiones de arquitectura y operar el sistema con un solo comando.
\end{enumerate}

% ============================================================================
\clearpage
\section{Marco teórico}
% ============================================================================

\subsection{Microservicios}

Un microservicio es un programa pequeño con una única responsabilidad, desplegable de forma
independiente y dueño exclusivo de sus datos. La comunicación entre servicios puede ser
\textbf{síncrona} (HTTP, cuando se necesita una respuesta) o \textbf{asíncrona} (eventos, cuando
solo se anuncia un hecho).

\subsection{Arquitectura orientada a eventos}

Los servicios publican \emph{hechos} (por ejemplo \pkg{shipment.created}) en un registro
distribuido. Kafka conserva los mensajes en orden por partición y permite que varios consumidores
lean el mismo flujo sin borrarlo. La clave de partición es el identificador del agregado, lo que
garantiza el orden por envío o por conductor.

\subsection{Patrón Outbox}

Escribir en la base y publicar en Kafka son dos operaciones distintas que no comparten
transacción. El patrón \emph{outbox} resuelve el problema escribiendo el evento en una tabla de
la misma base y en la misma transacción que el cambio de negocio; un proceso posterior lo publica
y solo entonces lo marca como enviado.

\subsection{Saga coreografiada}

Una operación que abarca varios servicios se modela como una saga: cada servicio reacciona a los
eventos de los demás y, si un paso falla, se ejecuta una \emph{compensación} (por ejemplo,
devolver el envío a \texttt{CREATED} cuando el repartidor abandona).

\subsection{Idempotencia}

Una operación es idempotente si ejecutarla dos veces produce el mismo efecto que ejecutarla una.
Se implementa con una clave enviada por el cliente (\pkg{Idempotency-Key}) y una reclamación
atómica en la base de datos.

\subsection{Geocercas e histéresis}

Una geocerca es un círculo en el mapa. Para evitar falsos positivos por el ruido del GPS
($\pm 30$\,m en entorno urbano) se exige que el conductor permanezca dentro del radio durante
varias lecturas consecutivas (histéresis).

\subsection{Autenticación con JWT y JWKS}

El servicio de autenticación firma tokens con RSA (RS256) y publica la llave pública en un
endpoint JWKS. El resto de servicios validan la firma localmente, sin compartir secretos y sin
consultar al emisor en cada petición.

% ============================================================================
\clearpage
\section{Arquitectura del sistema}
% ============================================================================

\subsection{Estilo arquitectónico}

La arquitectura es de \textbf{microservicios con base de datos por servicio} y comunicación mixta:
HTTP para consultas que requieren respuesta inmediata (por ejemplo, el cálculo de ETA) y Kafka
para la propagación de hechos de negocio.

\begin{nota}[Regla de oro]
Ningún servicio lee ni escribe la base de otro. Si necesita un dato, lo pide por HTTP o se
suscribe al evento que lo anuncia. Esta regla es la que permite cambiar un servicio sin tocar los
demás.
\end{nota}

\subsection{Infraestructura}

\begin{table*}[t]
\centering
\caption{Componentes de infraestructura}
\small
\begin{tabular}{@{}L{3.2cm}L{8.2cm}L{3cm}@{}}
\toprule
\textbf{Componente} & \textbf{Uso} & \textbf{Puerto host}\\
\midrule
Kafka (KRaft) & 5 tópicos de negocio + 4 DLQ, 3 particiones por tópico & 29092\\
PostgreSQL 16 & 4 bases independientes, una por servicio & 5433\\
Redis 7 & índice geográfico, última posición, read-model y caché & 6379\\
OSRM & rutas por calles con extract de Fusagasugá & 5000\\
Nominatim & geocodificación de direcciones (API pública de OSM) & 443\\
MailHog & servidor SMTP de desarrollo & 8025\\
Prometheus & recolección de métricas de los 7 servicios & 9090\\
Grafana & tablero provisionado & 3000\\
\bottomrule
\end{tabular}
\end{table*}

\subsection{Flujo de una petición}

\begin{enumerate}[leftmargin=1.2cm]
  \item El cliente envía \pkg{POST /api/v1/shipments} con su JWT y una \pkg{Idempotency-Key}.
  \item El gateway valida la firma del token y reenvía la petición a \pkg{shipment-service}.
  \item \pkg{shipment-service} guarda el envío, su historial y el evento en el outbox, todo en
        una transacción.
  \item El relay publica \pkg{shipment.created} en Kafka.
  \item \pkg{dispatch-service} consume el evento, busca candidatos en Redis y emite ofertas.
  \item El repartidor acepta; se crea la asignación y se publica \pkg{dispatch.assigned}.
  \item \pkg{shipment-service} cambia el estado a \texttt{ASSIGNED}.
\end{enumerate}

% ============================================================================
\clearpage
\section{Modelo de datos}
\label{ch:datos}
% ============================================================================

\subsection{Metodología}

El modelo se construyó en tres niveles: \textbf{conceptual} (entidades y relaciones),
\textbf{lógico} (tablas, claves y tipos) y \textbf{físico} (índices, restricciones y
migraciones Flyway). Cada servicio posee su propia base de datos; las tablas transversales
(\pkg{outbox}, \pkg{processed\_events}, \pkg{idempotency\_keys}, \pkg{shedlock}) se replican
solo en los servicios que las necesitan.

\subsection{Modelo entidad-relación (MER)}

La figura \ref{fig:mer} presenta el modelo conceptual. Se muestran las entidades principales con
sus atributos identificadores y las relaciones con sus cardinalidades.

\begin{figure*}[!t]
\centering
\resizebox{\textwidth}{!}{%
\begin{tikzpicture}[
  entity/.style={draw=ink, rounded corners=3pt, fill=paper, minimum width=3.9cm,
                 minimum height=1.1cm, align=center, font=\small, inner sep=4pt},
  attr/.style={font=\scriptsize\color{soft}},
  rel/.style={-{Latex}, thick, color=soft},
  card/.style={font=\scriptsize\bfseries, fill=white, inner sep=1.5pt}
]
  % columna izquierda / centro / derecha
  \node[entity] (role) at (0,10.8) {\textbf{ROLE}\\[1pt] \attr{name}};
  \node[entity] (refresh) at (6,10.8) {\textbf{REFRESH\_TOKEN}\\[1pt] \attr{id · token\_hash · expires\_at}};

  \node[entity] (user) at (0,6.8) {\textbf{USER}\\[1pt] \attr{id · email · password\_hash}};
  \node[entity] (shipment) at (6,6.8) {\textbf{SHIPMENT}\\[1pt] \attr{id · customer\_id · status}};
  \node[entity] (driver) at (12,6.8) {\textbf{DRIVER}\\[1pt] \attr{id · status · rating}};

  \node[entity] (state) at (0,3) {\textbf{DISPATCH\_STATE}\\[1pt] \attr{shipment\_id · phase · round}};
  \node[entity] (offer) at (6,3) {\textbf{OFFER}\\[1pt] \attr{id · shipment\_id · driver\_id}};
  \node[entity] (assign) at (12,3) {\textbf{ASSIGNMENT}\\[1pt] \attr{id · shipment\_id · driver\_id}};

  \node[entity] (outbox) at (0,-0.8) {\textbf{OUTBOX}\\[1pt] \attr{id · aggregate\_id · payload}};
  \node[entity] (notif) at (6,-0.8) {\textbf{NOTIFICATION}\\[1pt] \attr{id · shipment\_id · status}};
  \node[entity] (idem) at (12,-0.8) {\textbf{IDEMPOTENCY\_KEY}\\[1pt] \attr{key · request\_hash}};

  % relaciones
  \draw[rel] (user) -- node[card, left]{1} node[card, right, pos=0.88]{M} (role)
        node[midway, font=\scriptsize, fill=paper, inner sep=2pt, xshift=0.75cm]{user\_roles};
  \draw[rel] (user) -- node[card, above left]{1} node[card, below right]{N} (refresh);
  \draw[rel] (user) -- node[card, above left]{1} node[card, below right]{N} (shipment);
  \draw[rel] (shipment) -- node[card, left]{1} node[card, right, pos=0.85]{N} (offer);
  \draw[rel] (driver) -- node[card, above right]{1} node[card, below left]{N} (offer);
  \draw[rel] (shipment) -- node[card, above right]{1} node[card, below left]{N} (assign);
  \draw[rel] (driver) -- node[card, left]{1} node[card, right, pos=0.85]{N} (assign);
  \draw[rel] (shipment) -- node[card, above right]{1} node[card, below left]{1} (state);
  \draw[rel] (shipment) -- node[card, left]{1} node[card, right, pos=0.85]{N} (notif);
\end{tikzpicture}
}
\caption{Modelo entidad-relación conceptual. \pkg{OUTBOX}, \pkg{PROCESSED\_EVENT} e
\pkg{IDEMPOTENCY\_KEY} son tablas de infraestructura: no participan del dominio, pero sostienen
la consistencia (eventos, deduplicación e idempotencia).}
\label{fig:mer}
\end{figure*}

\subsection{Entidades y atributos}

\begin{table*}[!t]
\centering
\caption{Entidades del dominio y sus atributos}
\label{tab:entidades}
\begin{tabular}{@{}L{4.6cm}L{2.2cm}L{7.2cm}@{}}
\toprule
\textbf{Entidad} & \textbf{Clave} & \textbf{Atributos principales}\\
\midrule
USER & id & email (único), password\_hash, full\_name, enabled, created\_at\\
ROLE & name & ADMIN, CLIENT, DRIVER (catálogo)\\
REFRESH\_TOKEN & id & user\_id, token\_hash (único), expires\_at, revoked, created\_at\\
SHIPMENT & id & customer\_id, customer\_email, driver\_id, status, origen/destino (lat, lng,
dirección), pickup\_pin, dispatch\_rounds, version, created\_at, updated\_at\\
SHIPMENT\_STATUS\_HISTORY & id & shipment\_id, from\_status, to\_status, actor\_id, reason, occurred\_at\\
DRIVER & id & full\_name, vehicle\_type, status, current\_lat/lng, rating, version, updated\_at\\
OFFER & id & shipment\_id, driver\_id, status, expires\_at, round, distance\_meters, score, created\_at\\
ASSIGNMENT & id & shipment\_id, driver\_id, status, created\_at, ended\_at\\
DISPATCH\_STATE & shipment\_id & pickup\_lat/lng, phase, round, current\_offer\_id, updated\_at\\
NOTIFICATION & id & shipment\_id, channel, recipient, template, variables, status, attempts,
last\_error, next\_attempt\_at, created\_at, sent\_at\\
OUTBOX & id & aggregate\_id, event\_type, topic, payload, created\_at, published\_at, attempts\\
PROCESSED\_EVENT & event\_id + consumer & processed\_at\\
IDEMPOTENCY\_KEY & key & request\_hash, response\_status, response\_body, created\_at, expires\_at\\
\bottomrule
\end{tabular}
\end{table*}

\subsection{Relaciones y cardinalidades}

\begin{table*}[t]
\centering
\caption{Relaciones del modelo conceptual}
\small
\begin{tabular}{@{}L{4.6cm}L{1.6cm}L{6.6cm}@{}}
\toprule
\textbf{Relación} & \textbf{Card.} & \textbf{Regla de negocio}\\
\midrule
USER — REFRESH\_TOKEN & 1:N & un usuario puede tener varios refresh; cada uno se revoca al usarse\\
USER — ROLE & N:M & tabla puente \pkg{user\_roles}\\
USER — SHIPMENT & 1:N & un cliente crea muchos envíos\\
SHIPMENT — OFFER & 1:N & varias ofertas por ronda (hasta 3 simultáneas)\\
DRIVER — OFFER & 1:N & un conductor recibe ofertas de distintos envíos\\
SHIPMENT — ASSIGNMENT & 1:0..1 & como máximo una asignación activa por envío\\
DRIVER — ASSIGNMENT & 1:N & histórico de asignaciones del conductor\\
SHIPMENT — DISPATCH\_STATE & 1:1 & una fila de estado por envío\\
SHIPMENT — NOTIFICATION & 1:N & varios correos por envío (entrega, reintentos)\\
SHIPMENT — OUTBOX & 1:N & eventos generados por el envío\\
\bottomrule
\end{tabular}
\end{table*}

\subsection{Modelo relacional: cuatro bases de datos}

\begin{table*}[t]
\centering
\caption{Bases de datos y tablas}
\small
\begin{tabular}{@{}L{2.6cm}L{11.4cm}@{}}
\toprule
\textbf{Base} & \textbf{Tablas}\\
\midrule
\pkg{authdb} & users, user\_roles, refresh\_tokens\\
\pkg{shipmentdb} & shipments, shipment\_status\_history, outbox, idempotency\_keys,
processed\_events, shedlock\\
\pkg{dispatchdb} & drivers, offers, assignments, dispatch\_state, outbox, processed\_events, shedlock\\
\pkg{notifdb} & notifications, processed\_events, shedlock\\
\bottomrule
\end{tabular}
\end{table*}

\begin{alerta}[¿Y tracking?]
\pkg{tracking-service} \textbf{no tiene base de datos}. Su estado (última posición, contadores de
geocerca, read-model de envíos) vive en Redis y es reconstruible: si el servicio se reinicia, el
siguiente ping lo corrige. Tampoco usa outbox porque no hay transacción de base que proteger.
\end{alerta}

\subsection{Diccionario de datos}

A continuación se detalla cada tabla: columnas, tipos, nulabilidad y restricciones. Los tipos
corresponden a PostgreSQL 16.

\subsubsection{authdb}

\begin{table*}[!t]
\centering
\caption{\pkg{users}}

\begin{tabular}{@{}L{3.4cm}L{3.2cm}L{1.2cm}L{3.6cm}L{3.4cm}@{}}
\toprule
\textbf{Columna} & \textbf{Tipo} & \textbf{Nulo} & \textbf{Restricción} & \textbf{Descripción}\\
\midrule
id & UUID & no & PK & identificador\\
email & VARCHAR(255) & no & UNIQUE & correo normalizado a minúsculas\\
password\_hash & VARCHAR(255) & no & & BCrypt con coste 12\\
full\_name & VARCHAR(255) & no & & nombre visible\\
enabled & BOOLEAN & no & default TRUE & permite deshabilitar cuentas\\
created\_at & TIMESTAMPTZ & no & & alta del usuario\\
\bottomrule
\end{tabular}
\end{table*}

\begin{table*}[!t]
\centering
\caption{\pkg{user\_roles}}

\begin{tabular}{@{}L{3.4cm}L{3.2cm}L{1.2cm}L{3.6cm}L{3.4cm}@{}}
\toprule
\textbf{Columna} & \textbf{Tipo} & \textbf{Nulo} & \textbf{Restricción} & \textbf{Descripción}\\
\midrule
user\_id & UUID & no & PK, FK users & usuario\\
role & VARCHAR(32) & no & PK & ADMIN, CLIENT o DRIVER\\
\bottomrule
\end{tabular}
\end{table*}

\begin{table*}[!t]
\centering
\caption{\pkg{refresh\_tokens}}

\begin{tabular}{@{}L{3.4cm}L{3.2cm}L{1.2cm}L{3.6cm}L{3.4cm}@{}}
\toprule
\textbf{Columna} & \textbf{Tipo} & \textbf{Nulo} & \textbf{Restricción} & \textbf{Descripción}\\
\midrule
id & UUID & no & PK & identificador\\
user\_id & UUID & no & FK users & dueño del token\\
token\_hash & VARCHAR(64) & no & UNIQUE & SHA-256 del token (nunca el token en claro)\\
expires\_at & TIMESTAMPTZ & no & & caducidad (7 días)\\
revoked & BOOLEAN & no & default FALSE & se revoca al rotar o al cerrar sesión\\
created\_at & TIMESTAMPTZ & no & & emisión\\
\bottomrule
\end{tabular}
\end{table*}

\subsubsection{shipmentdb}

\begin{table*}[!t]
\centering
\caption{\pkg{shipments}}

\begin{tabular}{@{}L{3.2cm}L{3.2cm}L{1.1cm}L{3.4cm}L{3.9cm}@{}}
\toprule
\textbf{Columna} & \textbf{Tipo} & \textbf{Nulo} & \textbf{Restricción} & \textbf{Descripción}\\
\midrule
id & UUID & no & PK & identificador del envío\\
customer\_id & UUID & no & & usuario que lo solicita\\
customer\_email & VARCHAR(255) & sí & & correo para la notificación\\
driver\_id & UUID & sí & & conductor asignado\\
status & VARCHAR(24) & no & & CREATED, ASSIGNED, PICKED\_UP, IN\_TRANSIT, DELIVERED, FAILED,
CANCELLED, NO\_DRIVER\\
pickup\_lat / lng & DOUBLE & no & & coordenadas de recogida\\
dropoff\_lat / lng & DOUBLE & no & & coordenadas de entrega\\
pickup\_address & VARCHAR(300) & no & & dirección legible\\
dropoff\_address & VARCHAR(300) & no & & dirección legible\\
pickup\_pin & VARCHAR(4) & no & & PIN de recogida para el cliente\\
dispatch\_rounds & SMALLINT & no & default 0 & rondas de despacho consumidas\\
version & BIGINT & no & default 0 & bloqueo optimista de JPA\\
created\_at / updated\_at & TIMESTAMPTZ & no & & marcas de tiempo\\
\bottomrule
\end{tabular}
\end{table*}

\begin{table*}[!t]
\centering
\caption{\pkg{shipment\_status\_history}}

\begin{tabular}{@{}L{3.2cm}L{3.2cm}L{1.1cm}L{3.4cm}L{3.9cm}@{}}
\toprule
\textbf{Columna} & \textbf{Tipo} & \textbf{Nulo} & \textbf{Restricción} & \textbf{Descripción}\\
\midrule
id & BIGSERIAL & no & PK & secuencia\\
shipment\_id & UUID & no & FK shipments & envío\\
from\_status & VARCHAR(24) & sí & & estado anterior (nulo al crear)\\
to\_status & VARCHAR(24) & no & & estado nuevo\\
actor\_id & UUID & sí & & quién provocó el cambio\\
reason & VARCHAR(300) & sí & & texto libre\\
occurred\_at & TIMESTAMPTZ & no & & instante del cambio\\
\bottomrule
\end{tabular}
\end{table*}

\begin{table*}[!t]
\centering
\caption{\pkg{outbox}}

\begin{tabular}{@{}L{3.2cm}L{3.2cm}L{1.1cm}L{3.4cm}L{3.9cm}@{}}
\toprule
\textbf{Columna} & \textbf{Tipo} & \textbf{Nulo} & \textbf{Restricción} & \textbf{Descripción}\\
\midrule
id & UUID & no & PK & identificador del evento\\
aggregate\_id & UUID & no & & clave de partición en Kafka\\
event\_type & VARCHAR(64) & no & & tipo del evento\\
topic & VARCHAR(64) & no & & tópico destino\\
payload & JSONB & no & & envoltorio completo serializado\\
created\_at & TIMESTAMPTZ & no & & alta\\
published\_at & TIMESTAMPTZ & sí & & nulo hasta el ack de Kafka\\
attempts & SMALLINT & no & default 0 & reintentos\\
\bottomrule
\end{tabular}
\end{table*}

\begin{table*}[!t]
\centering
\caption{\pkg{idempotency\_keys}}

\begin{tabular}{@{}L{3.2cm}L{3.2cm}L{1.1cm}L{3.4cm}L{3.9cm}@{}}
\toprule
\textbf{Columna} & \textbf{Tipo} & \textbf{Nulo} & \textbf{Restricción} & \textbf{Descripción}\\
\midrule
key & VARCHAR(128) & no & PK & valor de \pkg{Idempotency-Key}\\
request\_hash & VARCHAR(64) & no & & SHA-256 del cuerpo\\
response\_status & INTEGER & sí & & 0 = en curso\\
response\_body & JSONB & sí & & respuesta guardada para repetirla\\
created\_at & TIMESTAMPTZ & no & & alta\\
expires\_at & TIMESTAMPTZ & no & & 24 horas\\
\bottomrule
\end{tabular}
\end{table*}

\begin{table*}[!t]
\centering
\caption{\pkg{processed\_events} (también en dispatchdb y notifdb)}

\begin{tabular}{@{}L{3.2cm}L{3.2cm}L{1.1cm}L{3.4cm}L{3.9cm}@{}}
\toprule
\textbf{Columna} & \textbf{Tipo} & \textbf{Nulo} & \textbf{Restricción} & \textbf{Descripción}\\
\midrule
event\_id & UUID & no & PK & identificador del evento consumido\\
consumer & VARCHAR(64) & no & PK & nombre del consumidor\\
processed\_at & TIMESTAMPTZ & no & & instante de procesamiento\\
\bottomrule
\end{tabular}
\end{table*}

\begin{table*}[!t]
\centering
\caption{\pkg{shedlock} (también en dispatchdb y notifdb)}

\begin{tabular}{@{}L{3.2cm}L{3.2cm}L{1.1cm}L{3.4cm}L{3.9cm}@{}}
\toprule
\textbf{Columna} & \textbf{Tipo} & \textbf{Nulo} & \textbf{Restricción} & \textbf{Descripción}\\
\midrule
name & VARCHAR(64) & no & PK & nombre del lock\\
lock\_until & TIMESTAMPTZ & no & & hasta cuándo está tomado\\
locked\_at & TIMESTAMPTZ & no & & cuándo se tomó\\
locked\_by & VARCHAR(255) & no & & instancia que lo tomó\\
\bottomrule
\end{tabular}
\end{table*}

\subsubsection{dispatchdb}

\begin{table*}[!t]
\centering
\caption{\pkg{drivers}}

\begin{tabular}{@{}L{3.2cm}L{3.2cm}L{1.1cm}L{3.4cm}L{3.9cm}@{}}
\toprule
\textbf{Columna} & \textbf{Tipo} & \textbf{Nulo} & \textbf{Restricción} & \textbf{Descripción}\\
\midrule
id & UUID & no & PK & mismo id que el usuario en authdb\\
full\_name & VARCHAR(120) & no & & nombre visible\\
vehicle\_type & VARCHAR(16) & no & & MOTO, BICI o CARRO\\
status & VARCHAR(16) & no & & AVAILABLE, BUSY u OFFLINE\\
current\_lat / lng & DOUBLE & sí & & última posición conocida\\
rating & DOUBLE & no & & calificación 0–5\\
version & BIGINT & no & default 0 & bloqueo optimista\\
updated\_at & TIMESTAMPTZ & no & & última actualización\\
\bottomrule
\end{tabular}
\end{table*}

\begin{table*}[!t]
\centering
\caption{\pkg{offers}}

\begin{tabular}{@{}L{3.2cm}L{3.2cm}L{1.1cm}L{3.4cm}L{3.9cm}@{}}
\toprule
\textbf{Columna} & \textbf{Tipo} & \textbf{Nulo} & \textbf{Restricción} & \textbf{Descripción}\\
\midrule
id & UUID & no & PK & identificador\\
shipment\_id & UUID & no & & envío ofertado\\
driver\_id & UUID & no & FK drivers & conductor candidato\\
status & VARCHAR(16) & no & & CREATED, ACCEPTED, REJECTED, EXPIRED, SUPERSEDED\\
expires\_at & TIMESTAMPTZ & no & & vencimiento (90 s)\\
round & SMALLINT & no & & ronda 1–3\\
distance\_meters & INTEGER & sí & & distancia al origen\\
score & DOUBLE & no & & puntuación 70\,\% cercanía + 30\,\% rating\\
version & BIGINT & no & default 0 & bloqueo optimista\\
created\_at & TIMESTAMPTZ & no & & emisión\\
\bottomrule
\end{tabular}
\end{table*}

\begin{table*}[!t]
\centering
\caption{\pkg{assignments}}

\begin{tabular}{@{}L{3.2cm}L{3.2cm}L{1.1cm}L{3.4cm}L{3.9cm}@{}}
\toprule
\textbf{Columna} & \textbf{Tipo} & \textbf{Nulo} & \textbf{Restricción} & \textbf{Descripción}\\
\midrule
id & UUID & no & PK & identificador\\
shipment\_id & UUID & no & & envío\\
driver\_id & UUID & no & FK drivers & conductor\\
status & VARCHAR(16) & no & & ACTIVE, CANCELLED, COMPLETED\\
created\_at & TIMESTAMPTZ & no & & inicio\\
ended\_at & TIMESTAMPTZ & sí & & fin\\
\bottomrule
\end{tabular}
\end{table*}

\begin{table*}[!t]
\centering
\caption{\pkg{dispatch\_state}}

\begin{tabular}{@{}L{3.2cm}L{3.2cm}L{1.1cm}L{3.4cm}L{3.9cm}@{}}
\toprule
\textbf{Columna} & \textbf{Tipo} & \textbf{Nulo} & \textbf{Restricción} & \textbf{Descripción}\\
\midrule
shipment\_id & UUID & no & PK & envío\\
pickup\_lat / lng & DOUBLE & no & & origen (copiado del evento)\\
phase & VARCHAR(16) & no & & SEARCHING, OFFERED, ASSIGNED, COMPLETED, EXHAUSTED, CANCELLED\\
round & SMALLINT & no & default 0 & ronda actual\\
current\_offer\_id & UUID & sí & & oferta destacada (informativo)\\
updated\_at & TIMESTAMPTZ & no & & último cambio\\
\bottomrule
\end{tabular}
\end{table*}

\subsubsection{notifdb}

\begin{table*}[!t]
\centering
\caption{\pkg{notifications}}

\begin{tabular}{@{}L{3.2cm}L{3.2cm}L{1.1cm}L{3.4cm}L{3.9cm}@{}}
\toprule
\textbf{Columna} & \textbf{Tipo} & \textbf{Nulo} & \textbf{Restricción} & \textbf{Descripción}\\
\midrule
id & UUID & no & PK & identificador\\
shipment\_id & UUID & sí & & nulo en notificaciones de prueba\\
channel & VARCHAR(16) & no & & EMAIL\\
recipient & VARCHAR(255) & no & & destinatario\\
template & VARCHAR(64) & no & & SHIPMENT\_DELIVERED, TEST\\
variables & JSONB & no & & datos de la plantilla\\
status & VARCHAR(16) & no & & PENDING, SENT, FAILED\\
attempts & SMALLINT & no & default 0 & intentos realizados\\
last\_error & VARCHAR(500) & sí & & último error\\
next\_attempt\_at & TIMESTAMPTZ & sí & & próximo reintento\\
created\_at & TIMESTAMPTZ & no & & alta\\
sent\_at & TIMESTAMPTZ & sí & & envío exitoso\\
\bottomrule
\end{tabular}
\end{table*}

\subsection{Índices y restricciones}

\begin{table*}[t]
\centering
\caption{Índices y restricciones relevantes}
\small
\begin{tabular}{@{}L{4.6cm}L{2.4cm}L{7.4cm}@{}}
\toprule
\textbf{Nombre} & \textbf{Tabla} & \textbf{Propósito}\\
\midrule
\pkg{uq\_offers\_pending\_driver} & offers & único parcial: una oferta pendiente por (envío, conductor)\\
\pkg{uq\_assignments\_active} & assignments & único parcial: una sola asignación activa por envío\\
\pkg{ix\_outbox\_pending} & outbox & parcial \pkg{WHERE published\_at IS NULL} para el relay\\
\pkg{ix\_notifications\_due} & notifications & parcial \pkg{WHERE status = 'PENDING'}\\
\pkg{ix\_shipments\_customer} & shipments & listado por cliente\\
\pkg{ix\_shipments\_status} & shipments & filtros por estado\\
\pkg{ix\_history\_shipment} & history & timeline ordenado\\
\pkg{ix\_drivers\_status} & drivers & disponibilidad\\
\pkg{ix\_offers\_expiry} & offers & barrido de vencidas\\
\pkg{ix\_refresh\_tokens\_user} & refresh\_tokens & revocación por usuario\\
\bottomrule
\end{tabular}
\end{table*}

\subsection{Modelo no relacional (Redis)}

\begin{table*}[t]
\centering
\caption{Claves de Redis}
\small
\begin{tabular}{@{}L{5.2cm}L{2cm}L{7.2cm}@{}}
\toprule
\textbf{Clave} & \textbf{Tipo} & \textbf{Contenido / TTL}\\
\midrule
\pkg{drivers:available} & GEO & conductores libres con posición (sin TTL)\\
\pkg{driver:\{id\}:last} & HASH & lat, lng, at (10 min)\\
\pkg{driver:\{id\}:shipment} & STRING & envío activo del conductor (12 h)\\
\pkg{shipment:\{id\}} & HASH & pickup, dropoff, driverId (12 h)\\
\pkg{geo:\{envio\}:\{cerca\}:hits} & STRING & contador de 3 pings (2 min)\\
\pkg{geo:\{envio\}:\{cerca\}:done} & STRING & geocerca confirmada (6 h)\\
\pkg{eta:\{lat,lng\}:\{lat,lng\}} & STRING & ruta cacheada (20 s)\\
\pkg{rl:gps:\{id\}} & STRING & límite de 3 lotes/s (1 s)\\
\bottomrule
\end{tabular}
\end{table*}

% ============================================================================
\clearpage
\section{Interfaz de programación (API)}
% ============================================================================

Todas las rutas se publican a través del gateway en \pkg{http://localhost:18080}. Salvo las
marcadas como públicas, requieren la cabecera \pkg{Authorization: Bearer <jwt>}.

\begin{table*}[!t]
\centering
\caption{Endpoints REST y WebSocket}
\label{tab:api}
\begin{tabular}{@{}L{6.4cm}L{2.2cm}L{1.6cm}L{5.2cm}@{}}
\toprule
\textbf{Método y ruta} & \textbf{Servicio} & \textbf{Auth} & \textbf{Descripción}\\
\midrule
POST /api/v1/auth/register & auth & público & crea un cliente\\
POST /api/v1/auth/login & auth & público & access + refresh token\\
POST /api/v1/auth/refresh & auth & público & rota el refresh\\
POST /api/v1/auth/logout & auth & público & revoca el refresh\\
GET /api/v1/auth/jwks & auth & público & llave pública RSA\\
GET /api/v1/auth/me & auth & jwt & perfil del token\\
POST /api/v1/shipments & shipment & jwt & crea envío (Idempotency-Key)\\
GET /api/v1/shipments & shipment & jwt & listado paginado\\
GET /api/v1/shipments/\{id\} & shipment & jwt & estado e historial\\
GET /api/v1/shipments/\{id\}/history & shipment & jwt & bitácora de estados\\
POST /api/v1/shipments/\{id\}/cancel & shipment & jwt & cancelación con reglas por rol\\
POST /api/v1/shipments/\{id\}/assign & shipment & ADMIN & asignación manual\\
GET /api/v1/shipments/summary & shipment & jwt & conteo por estado\\
POST /api/v1/shipments/\{id\}/confirm-pickup & shipment & jwt & recogida manual (sin GPS)\\
POST /api/v1/shipments/\{id\}/confirm-delivery & shipment & jwt & entrega manual (sin GPS)\\
POST /api/v1/drivers/\{id\}/location & dispatch & jwt & posición del conductor\\
POST /api/v1/drivers/\{id\}/availability & dispatch & jwt & disponibilidad\\
GET /api/v1/drivers/\{id\}/offers & dispatch & jwt & ofertas vigentes\\
POST /api/v1/offers/\{id\}/accept & dispatch & jwt & acepta (tres candados)\\
POST /api/v1/offers/\{id\}/reject & dispatch & jwt & rechaza y sigue\\
POST /api/v1/assignments/\{id\}/cancel & dispatch & jwt & compensación de la saga\\
GET /api/v1/dispatch/shipments/\{id\} & dispatch & jwt & fase, ronda y ofertas\\
POST /api/v1/tracking/drivers/\{id\}/location & tracking & jwt & lote GPS (202)\\
GET /api/v1/tracking/shipments/\{id\} & tracking & jwt & posición y ETA\\
GET /api/v1/tracking/shipments/\{id\}/trail & tracking & jwt & recorrido (200 puntos)\\
GET /api/v1/tracking/drivers/\{id\}/last & tracking & jwt & última posición\\
GET /api/v1/tracking/stats & tracking & jwt & envíos y conductores activos\\
GET /api/v1/eta & routing & jwt & distancia y duración\\
GET /api/v1/eta/route & routing & jwt & ruta con geometría\\
GET /api/v1/eta/distance & routing & jwt & solo distancia\\
POST /api/v1/eta/matrix & routing & jwt & matriz N×M\\
GET /api/v1/eta/nearest & routing & jwt & punto en la calle más cercana\\
GET /api/v1/eta/geocode & routing & jwt & Nominatim: dirección a coordenadas\\
GET /api/v1/notifications & notification & jwt & correos de un envío\\
GET /api/v1/notifications/\{id\} & notification & jwt & estado e intentos\\
POST /api/v1/notifications/\{id\}/retry & notification & ADMIN & reintento manual\\
POST /api/v1/notifications/\{id\}/cancel & notification & ADMIN & cancela una pendiente\\
GET /api/v1/notifications/failed & notification & ADMIN & últimas 50 fallidas\\
GET /api/v1/notifications/stats & notification & jwt & conteo por estado\\
POST /api/v1/notifications/test & notification & ADMIN & correo de prueba\\
WS /ws/tracking & tracking & jwt en query & posiciones en vivo\\
\bottomrule
\end{tabular}
\end{table*}

\subsection{Formato de errores}

Todos los errores usan \pkg{application/problem+json} (RFC 9457) con un campo \pkg{code}
estable. Ejemplos: \pkg{INVALID\_TRANSITION} (409), \pkg{OFFER\_ALREADY\_CLAIMED} (409),
\pkg{IDEMPOTENCY\_KEY\_REUSED} (409), \pkg{VALIDATION\_ERROR} (400), \pkg{RATE\_LIMITED} (429).

\subsection{APIs externas}

El proyecto consume dos servicios de OpenStreetMap, ambos aislados tras el
\pkg{routing-eta-service}:

\begin{itemize}[leftmargin=1.2cm]
  \item \textbf{OSRM.} Motor de rutas por calle. Se usa en cinco operaciones: ETA, ruta con
        geometr\'ia, distancia, matriz origen-destino y punto m\'as cercano. Corre en local con el
        extracto de Fusagasug\'a; si falla, el cortocircuito devuelve una l\'inea recta con la
        velocidad media configurada.
  \item \textbf{Nominatim.} Geocodificaci\'on inversa de direcciones a coordenadas. Exige un
        \pkg{User-Agent} identificable y limita a una petici\'on por segundo, por lo que la
        respuesta se cachea 24 horas en Redis.
\end{itemize}

El frontend consume adem\'as los \emph{tiles} de OpenStreetMap para dibujar el mapa. En total son
\textbf{tres dependencias externas}, dos de ellas consumidas por el backend.

% ============================================================================
\clearpage
\section{Eventos y mensajería}
% ============================================================================

\subsection{Tópicos}

\begin{table*}[t]
\centering
\caption{Tópicos de Kafka}
\small
\begin{tabular}{@{}L{5cm}L{2.2cm}L{2.4cm}L{5.2cm}@{}}
\toprule
\textbf{Tópico} & \textbf{Particiones} & \textbf{Retención} & \textbf{Productor}\\
\midrule
\pkg{shipment.events} & 3 & 7 días & shipment\\
\pkg{dispatch.events} & 3 & 7 días & dispatch\\
\pkg{tracking.gps.raw} & 6 & 1 hora & tracking (ingesta)\\
\pkg{tracking.events} & 3 & 7 días & tracking (geocercas)\\
\pkg{notification.commands} & 3 & 7 días & shipment\\
\pkg{*.dlq} & 1 & 14 días & consumidores\\
\bottomrule
\end{tabular}
\end{table*}

\subsection{Envoltorio de eventos}

\begin{verbatim}
{
  "eventId": "7f1c...",           // deduplicación
  "eventType": "shipment.assigned",
  "eventVersion": 1,
  "occurredAt": "2026-11-25T10:00:00Z",
  "aggregateId": "a1b2...",       // clave de partición
  "traceId": "4bf9...",
  "payload": { }
}
\end{verbatim}

\subsection{Garantías}

\begin{itemize}[leftmargin=1.2cm]
  \item \textbf{Orden por agregado:} la clave de partición es el id del envío o del conductor.
  \item \textbf{Al menos una vez:} los consumidores deduplican con \pkg{processed\_events}.
  \item \textbf{Sin pérdida en la publicación:} outbox + ack de Kafka antes de marcar enviado.
  \item \textbf{Mensajes venenosos:} reintentos con espera y luego DLQ.
\end{itemize}

% ============================================================================
\clearpage
\section{Seguridad}
% ============================================================================

\begin{table*}[t]
\centering
\caption{Matriz de permisos}
\small
\begin{tabular}{@{}L{5.6cm}L{2.2cm}L{2.2cm}L{2.2cm}L{2.2cm}@{}}
\toprule
\textbf{Operación} & \textbf{CLIENT} & \textbf{DRIVER} & \textbf{ADMIN} & \textbf{Público}\\
\midrule
Crear envío & sí (propio) & no & sí & no\\
Ver envío & solo el suyo & solo asignados & todos & no\\
Cancelar envío & hasta ASSIGNED & no & hasta IN\_TRANSIT & no\\
Aceptar oferta & no & solo la suya & sí & no\\
Reportar GPS & no & solo el suyo & sí & no\\
Listar notificaciones & no & no & sí & no\\
\bottomrule
\end{tabular}
\end{table*}

Los tokens son RS256, con 15 minutos de vida y refresh de 7 días almacenado como hash. El gateway
valida la firma contra el JWKS y cada servicio vuelve a validar el rol y la propiedad del recurso:
nunca se confía solo en el gateway.

% ============================================================================
\clearpage
\section{Pruebas y calidad}
% ============================================================================

\subsection{Estrategia}

\begin{itemize}[leftmargin=1.2cm]
  \item \textbf{Unitarias (Mockito):} reglas de dominio, servicios y controladores, sin
        infraestructura.
  \item \textbf{Integración (Testcontainers):} PostgreSQL, Redis y MailHog reales; se valida el
        esquema Flyway y la entrega real de correo.
  \item \textbf{Concurrencia:} 50 hilos aceptando la misma oferta.
  \item \textbf{Contrato:} validación de los eventos publicados contra los esquemas.
\end{itemize}

\begin{table*}[t]
\centering
\caption{Tests y cobertura por módulo}
\small
\begin{tabular}{@{}L{4.6cm}C{2cm}C{2.6cm}@{}}
\toprule
\textbf{Módulo} & \textbf{Tests} & \textbf{Cobertura de líneas}\\
\midrule
tracking-service & 70 & 97,9\,\%\\
dispatch-service & 84 & 100,0\,\%\\
shipment-service & 58 & 94,4\,\%\\
auth-service & 41 & 95,3\,\%\\
notification-service & 42 & 97,6\,\%\\
common & 32 & 100,0\,\%\\
routing-eta-service & 29 & 97,8\,\%\\
contracts & 4 & 100,0\,\%\\
api-gateway & 3 & 90,0\,\%\\
\midrule
\textbf{Total (9 m\'odulos)} & \textbf{367} & \textbf{97,7\,\%}\\
\bottomrule
\end{tabular}
\end{table*}

El umbral del 90\,\% se verifica en integraci\'on continua con \pkg{./scripts/coverage.sh 90}.
La m\'etrica excluye clases de arranque, configuraci\'on de Spring, DTOs y sembradores, y la
herramienta de demostraci\'on \pkg{tools/driver-simulator}; son clases sin reglas de negocio.

\subsection{Resultados de validaci\'on experimental}
\label{sec:resultados}

Todos los resultados fueron verificados en la m\'aquina de desarrollo (JDK 21, Docker Compose v2,
PostgreSQL 16, Kafka KRaft, Redis 7, OSRM con extracto de Fusagasug\'a de 4,3\,MB). La tabla
\ref{tab:resultados} resume la evidencia de extremo a extremo.

\begin{table*}[t]
\centering
\caption{Resultados experimentales verificados}
\label{tab:resultados}
\small
\begin{tabular}{@{}L{3.8cm}L{7.2cm}C{3.2cm}@{}}
\toprule
\textbf{Prueba} & \textbf{Procedimiento y evidencia} & \textbf{Resultado}\\
\midrule
Ciclo de vida completo & crear env\'io $\rightarrow$ oferta $\rightarrow$ aceptar $\rightarrow$
conducir por calles $\rightarrow$ entrega; WebSocket recibe posiciones; MailHog recibe correo &
\texttt{CREATED} $\rightarrow$ \texttt{DELIVERED} en $\sim$2 min \\
Asignaci\'on bajo contenci\'on & 50 hilos aceptan la misma oferta simult\'aneamente (test) y
2 peticiones HTTP paralelas (laboratorio) & 1 ganador, 49 conflictos (409); 1 asignaci\'on activa \\
Idempotencia & mismo \pkg{Idempotency-Key} + mismo cuerpo repetido; luego mismo key con cuerpo distinto &
mismo \texttt{id} devuelto; despu\'es 409 \pkg{IDEMPOTENCY\_KEY\_REUSED} \\
Geocercas con hist\'eresis & 3 pings consecutivos dentro de 100\,m del origen; alejarse $>$150\,m &
\texttt{PICKED\_UP} al tercer ping; luego \texttt{IN\_TRANSIT} \\
ETA real y cach\'e & \pkg{GET /api/v1/eta} Parque $\rightarrow$ Terminal; repetir llamada; apagar OSRM &
2\,529\,m / 273\,s \pkg{source=OSRM}; 2.a \pkg{cached:true}; sin OSRM \pkg{FALLBACK} en 43\,ms \\
Geocodificaci\'on externa & \pkg{GET /api/v1/eta/geocode?q=Fusagasug\'a}; repetir la misma consulta &
Nominatim responde 4,3449509 / $-74,3674249$ en 1,35\,s; 2.a consulta \pkg{cached:true} en 28\,ms \\
Carga y estr\'es & \pkg{load/ab.sh}: 2000 peticiones, 20 concurrentes, 7 endpoints &
p95 entre 16 y 46\,ms (l\'imite 500\,ms); 0 errores de red; quiebre hacia 400 concurrentes \\
Notificaci\'on confiable & entregar env\'io; consultar \pkg{/api/v1/notifications}; apagar MailHog y reintentar &
\texttt{SENT} con 1 intento; sin SMTP: 3 intentos (2\,s, 3\,s, 5\,s) $\rightarrow$ \texttt{FAILED} + DLQ \\
Saga de compensaci\'on & cancelar \pkg{assignment} como repartidor tras \texttt{ASSIGNED} &
env\'io vuelve a \texttt{CREATED} y se reoferta; conductor liberado \\
Observabilidad & Prometheus \pkg{/targets} + Grafana + \pkg{vite build} &
8/8 targets \texttt{UP}; dashboard operativo; frontend compila \\
\bottomrule
\end{tabular}
\end{table*}

Tres cifras sostienen la calificaci\'on del sistema: (i) \textbf{cero doble asignaci\'on} en 50
intentos concurrentes gracias a la triple barrera (Redis \pkg{SET NX}, \pkg{UPDATE \dots WHERE
status='CREATED'} e \'indice \'unico parcial); (ii) \textbf{cero eventos perdidos}, porque el
outbox retiene el evento si Kafka cae y el relay lo reintenta; y (iii) \textbf{cero env\'ios
duplicados} ante reintentos del cliente, por la reclamaci\'on at\'omica de la clave de
idempotencia.

% ============================================================================
\clearpage
\section{Despliegue y operación}
% ============================================================================

\subsection{Requisitos}

JDK 21, Maven 3.9+, Docker con Compose v2, Node 20+ y 8\,GB de RAM libres (16 recomendados).

\subsection{Arranque}

\begin{verbatim}
./scripts/start.sh        # infraestructura + 7 servicios + frontend
./scripts/tunel.sh        # URL HTTPS para el movil (Cloudflare)
./scripts/coverage.sh 70  # umbral de cobertura
./scripts/start.sh --stop # apagar todo
\end{verbatim}

\subsection{Observabilidad}

Prometheus recolecta \pkg{/actuator/prometheus} de los 7 servicios; Grafana muestra el tablero
\emph{EnRuta · plataforma} con peticiones por segundo, latencia p95, heap, aciertos de caché de
ETA e intentos de notificación.

% ============================================================================
\clearpage
\section{Decisiones de diseño}
% ============================================================================

\begin{table*}[t]
\centering
\caption{Resumen de los ADR}
\small
\begin{tabular}{@{}L{1cm}L{13cm}@{}}
\toprule
\textbf{ADR} & \textbf{Decisión}\\
\midrule
0001 & Kafka como broker de eventos (orden, replay y DLQ).\\
0002 & Un PostgreSQL con cuatro bases separadas en local.\\
0003 & Puertos no estándar (18080, 5433, 29092) por conflictos de la máquina.\\
0004 & Outbox con relay por polling y ShedLock.\\
0005 & JWT RS256 propio con JWKS.\\
0006 & Idempotencia con reclamación atómica (\pkg{INSERT ... ON CONFLICT}).\\
0007 & Doble (triple) barrera contra la doble asignación.\\
0008 & Saga coreografiada entre shipment y dispatch.\\
0009 & Tracking sin base de datos ni outbox.\\
0010 & WebSocket de una instancia y token en query string.\\
0011 & Ingesta GPS asíncrona con 202 y rate limit.\\
0012 & ETA con caché explícita, cortocircuito y plan B.\\
0013 & Notificaciones: la fila como outbox de salida.\\
\bottomrule
\end{tabular}
\end{table*}

% ============================================================================
\clearpage
\section{Conclusiones y trabajo futuro}
\label{ch:futuro}
% ============================================================================

El sistema demuestra que los problemas dif\'iciles de la \'ultima milla no est\'an en el CRUD sino en
la consistencia: asignar una sola vez, no perder eventos y sobrevivir a los fallos. Las tres
barreras de asignaci\'on, el outbox y la saga son las piezas que sostienen esa garant\'ia, y est\'an
respaldadas por 367 pruebas (97,7\,\%), prueba de concurrencia con un \'unico ganador, ciclo
completo en $\sim$2 min y correo entregado al primer intento.

\subsection{Trabajo futuro}

\begin{enumerate}[leftmargin=1.2cm]
  \item Multi-empresa (SaaS) con aislamiento por fila y API keys.
  \item Precio dinámico, pago y liquidación al repartidor.
  \item Aplicación móvil nativa con GPS en segundo plano.
  \item Trazas distribuidas completas (Tempo) y SLOs con alertas.
  \item Historial de rutas para analítica y ETA basado en datos históricos.
\end{enumerate}

% ============================================================================
\clearpage
\begin{thebibliography}{9}
\bibitem{hohpe} Gregor Hohpe, Bobby Woolf. \emph{Enterprise Integration Patterns}. Addison-Wesley, 2003.
\bibitem{kleppmann} Martin Kleppmann. \emph{Designing Data-Intensive Applications}. O'Reilly, 2017.
\bibitem{richardson} Chris Richardson. \emph{Microservices Patterns}. Manning, 2018.
\bibitem{spring} Spring Boot Reference Documentation, 3.4. \url{https://docs.spring.io/spring-boot/}.
\bibitem{kafka} Apache Kafka Documentation. \url{https://kafka.apache.org/documentation/}.
\bibitem{osrm} Project OSRM. \url{https://project-osrm.org/}.
\bibitem{rfc9457} M. Nottingham et al. \emph{Problem Details for HTTP APIs}, RFC 9457, 2023.
\end{thebibliography}

% ============================================================================
\clearpage
\appendices
\clearpage
\section{Glosario}

\begin{description}[leftmargin=3.2cm, style=nextline]
  \item[Microservicio] Programa con una responsabilidad y su propia base de datos.
  \item[Evento] Aviso de un hecho ocurrido, con identificador único.
  \item[Outbox] Tabla donde se escribe el evento en la misma transacción del cambio.
  \item[Idempotencia] Repetir una operación no cambia el resultado.
  \item[Saga] Secuencia de pasos con compensaciones.
  \item[DLQ] Cola de mensajes que agotaron sus reintentos.
  \item[Circuit breaker] Corta las llamadas a un servicio que falla y usa un plan B.
  \item[Geocerca] Círculo en el mapa que dispara una acción al entrar.
  \item[Histéresis] Exigir varias señales consecutivas para confirmar un cambio.
  \item[JWT / JWKS] Token firmado y llave pública para verificarlo.
  \item[Flyway] Migraciones versionadas de la base de datos.
  \item[Testcontainers] Contenedores levantados por las pruebas.
\end{description}

\clearpage
\section{Estructura del repositorio}

\begin{verbatim}
enruta/
|--  pom.xml                  parent Maven
|--  contracts/               eventos y nombres de topicos
|--  common/                  outbox, idempotencia, errores, geo
|--  services/
|  |--  api-gateway/
|  |--  auth-service/
|  |--  shipment-service/
|  |--  dispatch-service/
|  |--  tracking-service/
|  |--  routing-eta-service/
|  `--  notification-service/
|--  tools/driver-simulator/  conductores automaticos
|--  web/                     React + Leaflet (PWA)
|--  infra/                   docker-compose, Prometheus, Grafana, OSRM
|--  scripts/                 start, tunel, coverage, lab
|--  docs/                    ADRs, guias y PDFs
`--  load/k6/                 prueba de carga
\end{verbatim}


\end{document}
